\documentclass{article}
\usepackage{amsmath, amssymb, amsthm}
\usepackage{hyperref}
\usepackage{geometry}
\geometry{margin=1in}

\title{Erd\H{o}s Problem \#126}
\date{Last updated: 2025-08-31}
\author{Source: erdosproblems.com}

\begin{document}
\maketitle

\noindent\textbf{Status:} OPEN \quad \textbf{Prize: \$250}

\noindent\textbf{Tags:} number theory

\medskip

\noindent\textbf{Problem Statement:}

Let $f(n)$ be maximal such that if $A\subseteq\mathbb{N}$ has $\lvert A\rvert=n$ then $\prod_{a\neq b\in A}(a+b)$ has at least $f(n)$ distinct prime factors. Is it true that $f(n)/\log n\to\infty$?




Investigated by Erd\H{o}s and Tur\'{a}n \cite{ErTu34} (prompted by a question of L\'{a}z\'{a}r and Gr\"{u}nwald) in their first joint paper, where they proved that\[\log n \ll f(n) \ll n/\log n\](the upper bound is trivial, taking $A=\{1,\ldots,n\}$). Erd\H{o}s says that $f(n)=o(n/\log n)$ has never been proved, but perhaps never seriously attacked. This problem has been formalised in Lean as part of the Google DeepMind Formal Conjectures project .




References

[ErTu34] Erd\H{o}s, Paul and Turan, Paul, On a Problem in the Elementary Theory of Numbers . Amer. Math. Monthly (1934), 608-611.

\noindent\textbf{Related OEIS sequences:} possible


\bigskip
\noindent\small{Source: \url{https://www.erdosproblems.com/126}}
\end{document}
